\documentclass[11pt]{article}

\usepackage[margin=1in]{geometry}
\usepackage{amsmath,amssymb,amsthm}
\usepackage{graphicx}
\usepackage{hyperref}

\newtheorem{theorem}{Theorem}
\newtheorem{lemma}{Lemma}
\newtheorem{corollary}{Corollary}

\newcommand{\F}{\mathcal{F}}
\newcommand{\Lip}{\operatorname{Lip}}
\newcommand{\dist}{\operatorname{dist}}

\title{A Uniformly Discrete Coset-Quotient Subcase of a Question of Cuth--Doucha}
\author{arxiv\_automatic\_math run \texttt{fa\_banach\_001}}
\date{June 18, 2026}

\begin{document}
\maketitle

\section*{Source Question}

This packet concerns Question 5.5 from Marek Cuth and Michal Doucha,
\emph{Projections in Lipschitz-free spaces induced by group actions},
arXiv:2104.11519 \cite{CuthDoucha}.

\begin{center}
\includegraphics[width=\textwidth]{figures/open_problem_crop.png}
\end{center}

In transcription, the question asks:

\begin{quote}
Let $G$ be a metric group which is amenable and locally compact or SIN.
Let $H\subset G$ be its bounded subgroup. Is it true that
$\F(G/H) \stackrel{C}{\hookrightarrow} \F(G)$?
\end{quote}

\section*{Claimed Status}

Candidate partial result, likely valid pending human review. We prove the
subcase in which the coset quotient metric is uniformly discrete. The
general question remains open.

\section*{Definitions and Setup}

Let $(G,d)$ be a group equipped with a left-invariant metric and let $H$
be a subgroup. Write $e$ for the identity element and set
\[
  B=\sup_{h\in H} d(e,h).
\]
Thus $H$ is bounded exactly when $B<\infty$.

On the left-coset space $G/H$ consider
\[
  D(gH,fH)=\inf_{h,k\in H} d(gh,fk).
\]
We assume that $D$ is a metric, and that it is uniformly discrete: there
is a number $a>0$ such that $D(u,v)\ge a$ whenever $u\ne v$ in $G/H$.
The quotient map is denoted by $q:G\to G/H$.

\section*{Proof Intuition}

The obstruction in the noncompact bounded-subgroup problem is that the
usual averaging over $H$ need not land back in $\F(G)$; an invariant mean
naturally produces bidual objects. In the uniformly discrete quotient
case one can bypass averaging completely. Choose one representative in
each coset. Because every coset has diameter at most controlled by the
bounded subgroup $H$, this arbitrary transversal is coarse Lipschitz with
an additive error $2B$. Uniform discreteness of the quotient converts the
additive error into an honest Lipschitz bound. The universal property of
Lipschitz-free spaces then linearizes the quotient map and the section,
giving a bounded projection.

\begin{lemma}[Bounded fibers plus quotient separation give a Lipschitz section]
Let $(G,d)$, $H$, $B$, $D$, and $a$ be as above. Then every set-theoretic
section $s:G/H\to G$ of $q$ with $s(H)=e$ is Lipschitz, with
\[
  \operatorname{Lip}(s)\le 1+\frac{2B}{a}.
\]
\end{lemma}

\begin{proof}
Let $u,v\in G/H$. If $u=v$, there is nothing to prove. Suppose $u\ne v$.
Choose $\varepsilon>0$. By the definition of $D$, there are $h,k\in H$
such that
\[
  d(s(u)h,s(v)k)\le D(u,v)+\varepsilon.
\]
Since $d$ is left-invariant,
\[
  d(s(u),s(u)h)=d(e,h)\le B
  \quad\text{and}\quad
  d(s(v)k,s(v))=d(e,k^{-1})\le B.
\]
Therefore
\[
\begin{aligned}
  d(s(u),s(v))
  &\le d(s(u),s(u)h)+d(s(u)h,s(v)k)+d(s(v)k,s(v))\\
  &\le D(u,v)+2B+\varepsilon.
\end{aligned}
\]
Letting $\varepsilon\downarrow 0$ and using $D(u,v)\ge a$ gives
\[
  d(s(u),s(v))
  \le \left(1+\frac{2B}{a}\right)D(u,v).
\]
\end{proof}

\begin{theorem}
Under the hypotheses above, $\F(G/H,D)$ is linearly isomorphic to a
complemented subspace of $\F(G,d)$. More precisely, there is a projection
on $\F(G,d)$ onto a copy of $\F(G/H,D)$ with norm at most
$1+2B/a$.
\end{theorem}

\begin{proof}
The quotient map $q:(G,d)\to (G/H,D)$ is $1$-Lipschitz, because
$D(qg,qf)\le d(g,f)$ by taking $h=k=e$ in the infimum defining $D$.

Choose a section $s:G/H\to G$ with $s(H)=e$. By the lemma,
$s$ is $L$-Lipschitz with $L\le 1+2B/a$. Since $q\circ s$ is the identity
map on $G/H$, the induced linear maps on Lipschitz-free spaces satisfy
\[
  \widehat q\circ \widehat s = I_{\F(G/H)}.
\]
Consequently
\[
  P=\widehat s\circ \widehat q:\F(G,d)\to \F(G,d)
\]
is an idempotent operator. Its range is $\widehat s[\F(G/H,D)]$, and
the restriction of $\widehat q$ to this range is the inverse of
$\widehat s$. Thus the range is linearly isomorphic to $\F(G/H,D)$.
Moreover
\[
  \|P\|\le \|\widehat s\|\,\|\widehat q\|
  \le L\le 1+\frac{2B}{a}.
\]
\end{proof}

\section*{Verification and Scope}

The proof uses only the universal property of Lipschitz-free spaces and
the elementary section estimate above. No amenability, local compactness,
or SIN assumption is used. The price is the additional hypothesis that
the coset quotient metric is uniformly discrete.

This is a genuine partial result for Question 5.5, but it does not touch
the small-scale quotient case. If $D$ has nonzero distances tending to
zero, the estimate
\[
  d(s(u),s(v))\le D(u,v)+2B
\]
does not imply that $s$ is Lipschitz. That is the remaining obstruction
for this route.

\section*{Novelty and Literature Check}

Before promotion, the run checked the cheap indexes
\texttt{registry\_index.tsv}, \texttt{solutions/index.tsv},
\texttt{attempts/index.tsv}, and \texttt{proof\_gaps/index.tsv} for
arXiv:2104.11519 and the core phrases
\texttt{F(G/H)}, \texttt{bounded subgroup}, \texttt{Lip\_0(X/Y)}, and
\texttt{Lipschitz-free}. No exact duplicate packet was found.

The local parsed arXiv corpus found the exact question only in
arXiv:2104.11519. It also found the related compact-subgroup averaging
projection in Doucha--Kaufmann \cite{DouchaKaufmann}, which handles
compact $H$ by Haar averaging. The present packet is not a literature
identification of that result: it proves a different, elementary
uniformly discrete quotient subcase, including cases where no compact
averaging argument is available. Web searches for exact phrases around
\texttt{F(G/H) complemented F(G)} and \texttt{Lip\_0(X/Y) complemented}
did not reveal an exact later answer to the source question.

\section*{Human Review Recommendation}

Review the metric assumptions first: in particular, check that the
intended quotient metric $D$ is a genuine metric and that the quotient
has positive separation $a$. Under those assumptions, the proof reduces
to the displayed triangle inequality and the standard linearization
argument for Lipschitz-free spaces.

\begin{thebibliography}{9}

\bibitem{CuthDoucha}
M. Cuth and M. Doucha,
\emph{Projections in Lipschitz-free spaces induced by group actions},
arXiv:2104.11519, 2021.

\bibitem{DouchaKaufmann}
M. Doucha and P. L. Kaufmann,
\emph{Approximation properties in Lipschitz-free spaces over groups},
arXiv:2005.09785, 2020.

\bibitem{Weaver}
N. Weaver,
\emph{Lipschitz Algebras},
World Scientific, 1999.

\end{thebibliography}

\end{document}
